% ─────────────────────────────────────────────────────────────────
\section{Conclusion}
\label{sec:conclusion}
% ─────────────────────────────────────────────────────────────────

We presented a comprehensive study of protein family classification on the
SHREC\,25 benchmark, spanning the full arc from VTK mesh preprocessing to
hierarchical graph neural network training and detailed post-hoc analysis.

\paragraph{Pipeline.}
A two-stage pipeline converts raw VTK surface meshes into compact SSL graphs
(1\,000 nodes, 8D features) and then into geodesic Voronoi patch assignments
(64 patches), with all preprocessing implemented in pure Python and
parallelisable across all available CPU cores.

\paragraph{Architecture.}
The HierGIN model stacks node-level GIN message passing, attention-weighted
Voronoi patch pooling, patch-level GIN, and global max readout.
The two-level pooling hierarchy is the principal driver of the $+$15.8\,pp
improvement over flat GNN baselines.

\paragraph{Electrostatic potential is the key signal.}
Ablation experiments confirm that removing electrostatic potential degrades
macro-F1 by 11.8\,pp --- the largest single-feature impact observed.
NormalPotential adds a secondary gain, particularly for the hardest class
pairs.
Surface normals contribute a structural channel that improves recall on
geometrically distinct families.

\paragraph{Dataset ceiling.}
The pair (61,\,90) is near-inseparable under all current surface features
(SVM balanced accuracy 0.56, SW(Potential) = 0.088).
We estimate the ceiling with these features at macro-F1 $\approx$0.85--0.88;
our best model at 0.808 is within 4--7\,pp of this ceiling.
Closing the remaining gap requires sequence-level or volumetric information.

\paragraph{State of the art.}
HierGIN-14D (global norm) achieves macro-F1 $=$ \textbf{0.808}, accuracy
$=$ \textbf{0.865}, and balanced accuracy $=$ \textbf{0.812}, outperforming
both SHREC\,25 entries PointMamba (0.720) and Mamba3D (0.672).
The surface-topology-aware graph construction, combined with electrostatic
and geometric features with correct global normalisation, gives a structural
advantage over general point-cloud methods.

\paragraph{Enriched feature datasets.}
The 12D curvature-enriched dataset adds no separability beyond the 8D baseline
($\Delta\text{SVM} \leq +0.020$).
The 21D geometric-enriched dataset gains $+$0.085 SVM balanced accuracy for the
near-inseparable (61,\,90) pair.
Training HierGIN on a 14D subset (8D $+$ 6 geometric) with \emph{global}
normalisation (not per-protein z-scoring) achieves the best overall results:
accuracy 0.865, balanced accuracy 0.812, weighted-F1 0.862, macro-F1
\textbf{0.808}, outperforming the 8D baseline across all six metrics and
establishing a new state of the art on SHREC\,25.
The normalisation scheme is critical: per-protein z-scoring erases
cross-protein scale differences and degrades macro-F1 to 0.752.

\paragraph{Negative results.}
Two directions that did \emph{not} improve performance are worth recording.
Patch-level scalar re-injection --- feeding patch-mean raw electrostatic
values back between the two patch GNN layers --- consistently hurt performance
($-$1.8\,pp to $-$3.0\,pp macro-F1), ruling out ``signal loss through
pooling'' as an explanation for the remaining gap.
The failure is informative: the attention pooling already encodes the
relevant electrostatic information; explicit re-injection is redundant and
disrupts the learned representations.

\paragraph{Future work.}
Ongoing experiments evaluate a rotation-invariant node backbone (RI-GIN)
that encodes each edge message as a set of SO(3)-invariant scalars, removing
the dependence on global orientation.
Beyond this, promising directions include reducing the train-val overfitting
gap in the 14D model via regularisation or augmentation, contrastive
pre-training on full-resolution meshes, integration of protein language model
embeddings (e.g.\ ESM-2), and ensemble approaches that flag the
near-inseparable (61,\,90) pair for human review.
